\chapter{Run Summary and Monitoring}
\label{ch:run-summary}

\newthought{Every run ends with a summary} and can be watched while it runs. Together they tell you
whether a scan went well---how many points succeeded, how fast, and how hard the machine worked.

\section{The run summary files}
\label{sec:rs-files}

When a scan finishes, \Jarvis\ writes the same summary in three forms into the run's output
directory, and also prints it to the terminal and appends it to the main log:

\begin{description}[style=nextline, leftmargin=1.2em]
  \item[\file{run\_summary.json}] a machine-readable flat dictionary.
  \item[\file{run\_summary.csv}] a single-row table with the same fields.
  \item[\file{run\_summary.txt}] the human-readable terminal tables.
\end{description}

\section{What it reports}
\label{sec:rs-metrics}

The summary covers timing, throughput, parallelism, and resources
(Table~\ref{tab:rs-metrics}).

\begin{table}[ht]
  \footnotesize
  \begin{tabular}{@{}lp{0.58\linewidth}@{}}
    \toprule
    \textbf{Field} & \textbf{Meaning} \\
    \midrule
    \code{wall\_time\_sec}            & total elapsed time of the run \\
    \code{total\_points\_submitted}   & points handed to the factory \\
    \code{total\_points\_finished}    & points completed successfully \\
    \code{total\_points\_failed}      & points that failed \\
    \code{success\_rate}             & finished $/$ submitted \\
    \code{throughput\_points\_per\_min} & finished points per minute \\
    \code{avg/median\_point\_eval\_sec} & per-point evaluation time \\
    \code{configured/peak/mean\_active\_workers} & parallelism: set, peak, time-weighted mean \\
    \code{time\_in\_external\_tools\_sec} & cumulative time inside calculator commands \\
    \code{cpu\_percent\_mean}, \code{memory\_rss\_mb\_peak} & best-effort resource use (needs \code{psutil}) \\
    \bottomrule
  \end{tabular}
  \caption{Selected run-summary fields.}
  \label{tab:rs-metrics}
\end{table}

\begin{jwarn}
The execution-breakdown totals are \emph{cumulative across points}, not wall-clock. Because many
points run in parallel, the summed ``external tool'' and ``internal workflow'' times can exceed
\code{wall\_time\_sec}---they are not a measure of wasted time. The report says so directly:
``tracked cumulative timing across completed points; not equal to wall-clock time.''
\end{jwarn}

\section{The shutdown tables}
\label{sec:rs-shutdown}

At the end of a run you will also see two short tables in the terminal. The \textsc{hdf5} writer
reports how many records it enqueued and flushed; the factory reports the point tallies and the
average rate:

\begin{plaincode}
Factory shutdown summary
  submitted   10028
  ok          10028
  failed          0
  total_time  00:03:05
  avg_rate    53.93 points/s
\end{plaincode}

\noindent A clean run shows \code{failed} at (or near) zero and \code{ok} equal to
\code{submitted}.

\section{Watching a live run}
\label{sec:rs-monitor}

For a long scan, open a real-time resource monitor in a second terminal with \cli{-{}-monitor}
(Chapter~\ref{ch:cli}):

\begin{shellbox}
Jarvis bin/my_scan.yaml --monitor
\end{shellbox}

\noindent It refreshes every few seconds and shows \textsc{cpu}/memory use for the running scan;
press \code{q} to quit. Because it attaches to the live process, start it while the scan is
running.

\section{Benchmarking throughput}
\label{sec:rs-benchmark}

To measure raw throughput---for tuning workers or comparing machines---use \cli{-{}-benchmark},
which runs in a throughput mode for a fixed number of seconds and reports points-per-second:

\begin{shellbox}
Jarvis bin/my_scan.yaml --benchmark 60
\end{shellbox}

\noindent Performance tuning based on these numbers is covered in Chapter~\ref{ch:performance}.

\section{Summary}
\label{sec:rs-summary}

Read \file{run\_summary.txt} after a run for the headline numbers (success rate, throughput,
parallelism), remember the breakdown totals are cumulative not wall-clock, watch long runs with
\cli{-{}-monitor}, and measure raw speed with \cli{-{}-benchmark}.
